% Figura 3: fluxo da comparação estatística pareada entre dois algoritmos.
% Somente estrutura; nenhum valor numérico.
\documentclass[border=3pt]{standalone}
% Estilo comum das figuras esquemáticas (TikZ).
% Paleta Okabe--Ito (segura para daltonismo). A cor nunca é o único código:
% todo bloco de treino/teste também traz o rótulo textual.
\usepackage[T1]{fontenc}
\usepackage[utf8]{inputenc}
\usepackage[brazilian]{babel}
\usepackage{lmodern}
\usepackage{amsmath,amssymb}
\usepackage{tikz}
\usetikzlibrary{positioning,calc,decorations.pathreplacing,arrows.meta,fit,backgrounds}

\definecolor{cTreino}{HTML}{0072B2}   % azul
\definecolor{cTeste}{HTML}{E69F00}    % laranja
\definecolor{cModA}{HTML}{009E73}     % verde-azulado
\definecolor{cModB}{HTML}{CC79A7}     % púrpura-avermelhado
\definecolor{cAlerta}{HTML}{D55E00}   % vermelhão
\definecolor{cTinta}{HTML}{222222}
\definecolor{cCinza}{HTML}{6B6B6B}

\tikzset{
  every picture/.style={font=\sffamily\small, line cap=round, line join=round, text=cTinta},
  treino/.style={draw=cTreino!80!black, fill=cTreino!22, line width=0.5pt},
  teste/.style={draw=cTeste!70!black, fill=cTeste!45, line width=0.5pt},
  caixa/.style={draw=cTinta, rounded corners=2pt, line width=0.6pt, align=flush center,
                inner sep=4pt, fill=white},
  caixaA/.style={caixa, draw=cModA!80!black, fill=cModA!14},
  caixaB/.style={caixa, draw=cModB!80!black, fill=cModB!18},
  caixaT/.style={caixa, draw=cTreino!80!black, fill=cTreino!12},
  caixaTe/.style={caixa, draw=cTeste!70!black, fill=cTeste!25},
  caixaS/.style={caixa, draw=cCinza, fill=black!4},
  seta/.style={-{Stealth[length=2.4mm]}, line width=0.6pt, draw=cTinta},
  nota/.style={font=\sffamily\footnotesize, text=cCinza, align=flush left},
}

\begin{document}
\begin{tikzpicture}[x=1cm,y=1cm]
  % Dados e partições
  \node[caixaS, text width=10.4cm] (dados) at (8,0)
    {\textbf{Conjunto rotulado $\mathcal{D}$} ($n$ exemplos; grupos de dependência, se houver)};
  \node[caixaS, text width=10.4cm] (part) at (8,-1.35)
    {\textbf{Gerador de partições} com semente fixa, estratificado (e por grupos)\\
     $\Rightarrow$ manifesto de partições versionado e com \emph{hash}};
  \draw[seta] (dados) -- (part);
  \node[nota, anchor=west, text width=2.9cm] at (13.45,-1.35) {as mesmas partições servem aos algoritmos $A$ e $B$};

  % Caixa de cada partição
  \node[caixaT, text width=6.0cm] (treino) at (3.6,-3.45)
    {\textbf{Treino} $\mathcal{D}^{\mathrm{tr}}_{rj}$ ($n_{\mathrm{tr}}$ exemplos)};
  \node[caixaTe, text width=4.6cm] (teste) at (12.6,-3.45)
    {\textbf{Teste} $\mathcal{D}^{\mathrm{te}}_{rj}$ ($n_{\mathrm{te}}$ exemplos)};
  \node[caixaT, text width=6.0cm, align=flush left] (selecao) at (3.6,-5.25)
    {\textbf{Seleção interna, só no treino}\\[1pt]
     \footnotesize hiperparâmetros por validação cruzada interna;\\
     épocas (\emph{early stopping}) em validação interna;\\
     limiar e calibração por predições fora da dobra};
  \node[caixaA, text width=2.45cm] (A) at (2.1,-7.25) {\textbf{Algoritmo $A$}\\\footnotesize ajuste final};
  \node[caixaB, text width=2.45cm] (B) at (5.1,-7.25) {\textbf{Algoritmo $B$}\\\footnotesize ajuste final};
  \node[caixa, text width=4.6cm] (aval) at (12.6,-6.3)
    {\textbf{Avaliação no mesmo teste}\\ $M^{A}_{rj}$ e $M^{B}_{rj}$ (p.\,ex., MCC)};
  \node[caixa, text width=4.6cm] (dif) at (12.6,-7.85)
    {\textbf{Diferença pareada}\\ $d_{rj}=M^{A}_{rj}-M^{B}_{rj}$};
  \draw[seta] (part.south -| treino.north) -- (treino.north);
  \draw[seta] (part.south -| teste.north) -- (teste.north);
  \draw[seta] (treino) -- (selecao);
  \draw[seta] (selecao.south -| A.north) -- (A.north);
  \draw[seta] (selecao.south -| B.north) -- (B.north);
  \draw[seta] (teste) -- (aval);
  \draw[seta] (aval) -- (dif);
  \draw[seta, draw=cModA!80!black] (A.south) |- (8.9,-8.2) |- ($(aval.west)+(0,-0.15)$);
  \draw[seta, draw=cModB!80!black] (B.east) -- (8.6,-7.25) |- ($(aval.west)+(0,0.15)$);
  % Barreira: o teste não retorna à seleção
  \draw[cAlerta, line width=1.1pt, dash pattern=on 4pt off 2.5pt] (7.55,-3.0) -- (7.55,-6.0);
  \node[font=\sffamily\footnotesize, text=cAlerta, align=flush center, fill=white, inner sep=1pt] at (7.55,-4.4)
    {o teste\\não entra\\na seleção};
  % Contêiner
  \begin{scope}[on background layer]
    \node[draw=cCinza, dashed, rounded corners=4pt, fit=(treino)(teste)(selecao)(A)(B)(dif), inner sep=7pt,
          label={[anchor=south, font=\sffamily\small\bfseries, text=cCinza]north:para cada partição $(r,j)$}] (cont) {};
  \end{scope}

  % Coleta das diferenças e testes
  \node[caixaS, text width=8.6cm] (col) at (8,-10.1)
    {\textbf{Conjunto das diferenças pareadas} $\{d_{rj}\}$ de um único conjunto de dados};
  \draw[seta] (dif.south) |- ($(col.north)+(1.5,0.35)$) -- ($(col.north)+(1.5,0)$);
  \node[caixa, text width=6.6cm, align=flush left] (f) at (3.9,-12.45)
    {\textbf{5$\times$2cv} ($R=5$, $K=2$): 10 diferenças, 5 variâncias $s_r^2$\\
     $\Rightarrow$ $F$ combinado de Alpaydin, referência $F_{10,5}$\\[1pt]
     \footnotesize relatar: $p$; $\bar d$ e as 10 diferenças (descritivo)};
  \node[caixa, text width=6.6cm, align=flush left] (t) at (12.1,-12.45)
    {\textbf{$R\times K$-fold} (p.\,ex., $10\times10$): $RK$ diferenças\\
     $\Rightarrow$ $t$ reamostrado corrigido, referência $t_{RK-1}$\\[1pt]
     \footnotesize relatar: $\bar d$, IC corrigido e $p$};
  \draw[seta] (col.south -| f.north) -- (f.north);
  \draw[seta] (col.south -| t.north) -- (t.north);
  \node[nota, align=flush center] at (8,-11.15) {um teste por esquema;\\não intercambiáveis};
  \node[caixaS, dashed, text width=14.6cm] (multi) at (8,-14.45)
    {\textbf{Vários conjuntos de dados:} um efeito por conjunto (ou por grupo de conjuntos dependentes) $\Rightarrow$
     Wilcoxon (dois algoritmos) ou Friedman com comparações pós-teste, com correção de Holm};
  \draw[seta] (f.south) -- (f.south |- multi.north);
  \draw[seta] (t.south) -- (t.south |- multi.north);
\end{tikzpicture}
\end{document}
